% --- Секция: Исследование производительности и соответствие ТЗ ---

\section{Результаты исследования производительности}

В данном разделе представлены результаты тестирования разработанного красно-черного дерева (\texttt{RBTree}) в сравнении со стандартным контейнером \texttt{std::map} и детальный анализ использования ресурсов.

\subsection{Анализ потребления оперативной памяти (Valgrind)}

Для проверки эффективности техники \textit{Pointer Tagging} и контроля утечек памяти использовалась утилита \texttt{Valgrind Memcheck}. Сводные данные представлены ниже.

\textbf{RBTree}
\begin{verbatim}
==787525== HEAP SUMMARY:
==787525==     in use at exit: 122,880 bytes in 6 blocks
==787525==   total heap usage: 280,970 allocs, 280,964 frees, 20,425,944 bytes allocated
...
786090 LEAK SUMMARY:
786090    definitely lost: 0 bytes in 0 blocks
786090    indirectly lost: 0 bytes in 0 blocks
786090      possibly lost: 0 bytes in 0 blocks
786090    still reachable: 122,880 bytes in 6 blocks
786090         suppressed: 0 bytes in 0 blocks
786090
786090 ERROR SUMMARY: 0 errors from 0 contexts (suppressed: 0 from 0)
\end{verbatim}


\textbf{std::map}
\begin{verbatim}
==786090== HEAP SUMMARY:
==786090==     in use at exit: 122,880 bytes in 6 blocks
==786090==   total heap usage: 280,971 allocs, 280,965 frees, 18,178,304 bytes allocated
...
==787525== LEAK SUMMARY:
==787525==    definitely lost: 0 bytes in 0 blocks
==787525==    indirectly lost: 0 bytes in 0 blocks
==787525==      possibly lost: 0 bytes in 0 blocks
==787525==    still reachable: 122,880 bytes in 6 blocks
==787525==         suppressed: 0 bytes in 0 blocks
==787525== 
==787525== For lists of detected and suppressed errors, rerun with: -s
==787525== ERROR SUMMARY: 0 errors from 0 contexts (suppressed: 0 from 0)
\end{verbatim}

\begin{table}[h!]
\centering
\renewcommand{\arraystretch}{1.2}
\begin{tabular}{@{}lccc@{}}
\toprule
\textbf{Показатель} & \textbf{RBTree (Custom)} & \textbf{std::map (STL)} & \textbf{Разница} \\ \midrule
Total heap usage & 18,178,304 bytes & 20,425,944 bytes & \textbf{$-11.0\%$} \\
Размер узла (Node) & 64 bytes & 72 bytes & \textbf{$-11.0\%$} \\ \bottomrule
\end{tabular}
\end{table}

\textbf{Вывод по памяти:} Кастомная реализация потребляет на 2.14 МБ меньше памяти. Отсутствие \texttt{definitely lost} подтверждает корректную работу деструкторов. Наличие \texttt{still reachable} (122.8 КБ) относится к внутренним буферам стандартной библиотеки ввода-вывода и не является утечкой программы.

\subsection{Исследование скорости выполнения}

Замеры времени проводились на операциях вставки, поиска и удаления ключей типа \texttt{std::string}. Для каждой реализации было проведено по 10 тестов по 1 000 000 и 2 000 000 операций и посчитано среднее значение.
Тесты представляли из себя набор из случайных команд вставки, поиска и удаления в одинаковом количестве.

\begin{table}[h!]
\centering
\renewcommand{\arraystretch}{1.2}
\begin{tabular}{@{}lccc@{}}
\toprule
\textbf{Кол-во операций} & \textbf{RBTree (us)} & \textbf{std::map (us)} & \textbf{Разница} \\ \midrule
1,000,000 & 2,213,091 & 2,253,288 & $+1.8\%$ (быстрее) \\
2,000,000 & 4,482,729 & 4,726,862 & $+5.2\%$ (быстрее) \\ \bottomrule
\end{tabular}
\end{table}

\subsection{Анализ производительности (Callgrind)}

Для детального анализа производительности был использован инструмент \texttt{Valgrind Callgrind}, позволяющий измерить количество выполненных инструкций и выявить наиболее затратные участки программы.

\begin{verbatim}
summary: 9,958,524,374 Ir
\end{verbatim}

Общее количество выполненных инструкций составляет $\approx 9.96 \cdot 10^9$.

\subsubsection*{Неочевидные источники затрат}

Профилирование выявило, что значительная часть времени выполнения связана не с алгоритмами балансировки дерева, а с сопутствующими операциями:

\begin{itemize}
    \item \textbf{Посимвольное сравнение строк:}
    \begin{itemize}
        \item функции сравнения внутри \texttt{std::string} 
        (например, \texttt{char\_traits::compare})
        \item суммарно $\approx 10^8$ инструкций
    \end{itemize}

    \item \textbf{Оператор сравнения строк:}
    \begin{itemize}
        \item \texttt{std::operator==}
        \item $\approx 1{,}000{,}000$ вызовов
        \item $\approx 1.08 \cdot 10^8$ инструкций
    \end{itemize}

    \item \textbf{Внутренняя реализация поиска:}
    \begin{itemize}
        \item функция \texttt{search\_impl}
        \item значительная доля времени уходит на последовательные переходы по узлам и сравнение ключей
    \end{itemize}

    \item \textbf{Основной цикл программы:}
    \begin{itemize}
        \item функция \texttt{main}
        \item аккумулирует значительную часть инструкций за счёт генерации операций и работы с данными
    \end{itemize}
\end{itemize}

\subsubsection*{Вывод по профилированию}

Анализ показывает, что:

\begin{itemize}
    \item Существенная доля затрат приходится на \textbf{посимвольные операции над строками}, а не на структуру дерева.
    \item Даже при асимптотике $O(\log n)$ фактическая стоимость операций определяется длиной ключей.
    \item Функция \texttt{search\_impl} демонстрирует, что узким местом является не навигация по дереву, а сравнение ключей на каждом шаге.
\end{itemize}

\textbf{Практический вывод:} при использовании строковых ключей производительность определяется не только структурой данных, но и стоимостью операций сравнения, что важно учитывать при проектировании высоконагруженных систем.

\section{Дневник выполнения работы}

\begin{itemize}
    \item \textbf{Этап 1:} Реализация базовой структуры дерева. Реализация логики поворотов, балансировки при вставке и удалении.
    \item \textbf{Этап 2:} Оптимизация \texttt{Node} с помощью \textit{Pointer Tagging}.
    \item \textbf{Этап 3:} Рефакторинг, исправление недочётов и мелких ошибок. Проверка на утечки памяти с помощью \texttt{Valgrind}.
    \item \textbf{Этап 4:} Сравнение с \texttt{std::map} по скорости работы, по затрате памяти с помощью \texttt{Valgrind}. Анализ количества вызовов функций через \texttt{Callgrind} и оформление отчета.
\end{itemize}

\section{Выявленные недочёты и исправления}

В ходе разработки были устранены следующие проблемы:
\begin{enumerate}
    \item \textbf{Оверхед памяти:} Изначально цвет хранился как \texttt{enum}, что увеличивало размер узла из-за \textit{padding}. Внедрен \textbf{Pointer Tagging}, интегрировавший цвет в младший бит указателя на родителя.
    \item \textbf{Копирование данных:} Выявлено избыточное копирование строк при вставке. Добавлена перегрузка методов с \texttt{std::string\&\&}.
\end{enumerate}
